\section{引言}

卫星海洋遥感大数据与自主水下航行器的结合，正在将海洋监测从预设航线走航转变为先验地图引导的现场自适应巡测\cite{fossum2019}。多年累积的遥感产品提供了大尺度环境背景与演变统计，航行器则深入云层遮挡与遥感平滑难以企及的水体采集高分辨率样本。然而机动平台受制于有限的机载电池能量、嵌入式计算速率与固定作业时间窗口\cite{xiao2020thinking, magnino2026metareasoning}。航行器获取新观测后，既可调用机载高功耗芯片开展密集航线推演以优化后续路径，也可维持低功耗处理器后台规划而将电能与时间用于巡航采样或原位复测\cite{horvitz1987reasoning, russell1991right, shen2025thinking}。由于内部计算与外部观测共享同一块电池和同一任务窗口，何时继续计算、何时继续观测，是大数据驱动自主海洋观测的核心控制问题\cite{hunt2021formalizing}。

东海有害赤潮应急巡测集中体现了先验底图失真、水下通信受阻与机载资源受限带来的决策困难\cite{fossum2019, lou2014diurnal}。2023年中国海洋灾害公报记录全年赤潮46次、累计面积约$1500\text{ km}^2$，其中有毒有害赤潮29次\cite{mnr2024bulletin}。浙江台州、温州与宁波外海晚春至初夏常暴发东海原甲藻赤潮\cite{zhao2025distinguishing}，表层斑块在潮汐作用下往复漂移演变\cite{lou2014diurnal}。应急指挥中心通常向$80\text{ km}\times 80\text{ km}$海域派出一台百米级航行器开展24小时巡测，要求尽早定位叶绿素最浓水体并随时响应位置查询\cite{remus100}。受多星产品合成延迟、云区经验正交函数插补平滑\cite{beckers2003dineof}、水色反演偏差\cite{ye2021cocts}及船只航渡影响，航行器下水前装入的常是3天前旧图；加之水下几乎无法实时通信，航线重规划与资源调配均须由机载计算机独立完成\cite{jetsonorinnx}。

传统海洋巡测规则与预设时间分配策略难以在缺乏历史任务数据时适配不同机载硬件。一方面，割草机式全覆盖\cite{galceran2013coverage}、贪心追踪、上置信界采样\cite{srinivas2012gpucb}与汤普森采样\cite{russo2018thompson}等启发式规则将每步计算开销视为常数，忽略了多步前瞻规划对规避局部次优峰值的作用\cite{hollinger2014sampling}。另一方面，在多步规划器上采用固定时间分配策略（如仅靠后台规划的“一直走”或周期性开启高功耗芯片的规则），其表现高度依赖处理器速度与芯片功耗。当更换计算模块或电池配置时，需依靠该型硬件在大量历史赤潮任务中的实测表现才能选出合适策略，而新平台首次投入应急巡测时并无历史档案。

从历史任务中训练策略选择器同样受制于高昂采集成本与剧烈的任务间环境波动。近期机器人元推理研究尝试由历史任务训练回归或分类模型，根据出发前环境特征选择计算预算或用时策略\cite{sung2021quit, magnino2026metareasoning}。然而赤潮现场的初值格局、峰值位置与入水起点千差万别，任务窗口间的天然遗憾波动远大于优良策略间的平均性能差距，基于旧图静态特征训练的选择器极易被背景噪声淹没，且需要上百次历史任务作为训练样本。

经典在线单步元推理虽无需历史任务训练，但因轻信规划器在搜索样本上的自评得分而易陷入过度规划\cite{horvitz1987reasoning, russell1991right, cuadron2025overthinking}。该方法在每个周期比较“多算一步的得分增量”与“耗电及延误代价”来决定是否开启高功耗芯片或停下复测\cite{lin2015metareasoning, elboher2023formal}。由于航行器只能在一组由旧图生成的想象海洋上搜索航线，当规划器在同一批有限样本上既搜索航路点又评估预期收益时，优化过程会迎合特定样本中的随机正涨落，产生“优化器诅咒”\cite{smith2006curse, vanhasselt2010double}。即便额外推演在真实海洋中已不再改善航线，内部自评得分仍随推演次数持续虚高，误导控制器频繁开启高功耗芯片并中途停船\cite{jiang2015horizon, cuadron2025overthinking}。

针对上述问题，本文的核心思想是：航行器不应依据规划器自身的搜索得分决定是否继续计算，而应利用多年卫星大数据标定误差的留出想象海洋作为独立标尺，在未参与搜索的想象海洋上评判计算与观测结果（图\ref{fig:overview}）。本文用八年东海卫星叶绿素数据标定滞后底图误差，围绕旧图生成48个随现场采样经集合调整卡尔曼滤波\cite{anderson2001eakf, gaspari1999correlation}同步更新的想象海洋，划分为32个搜索用海洋与16个留出海洋。基于“搜索与评判分离”准则，本文提出两种无需历史任务的控制方式：下水前，由母船或岸基计算机以4个留出海洋为假想真值，按机载算力与功耗将五种策略各完整预演一遍24小时任务并择优执行；航行中，在单步元推理中改由16个留出海洋评估计算与采样收益并配合两倍标准误切换判据，在离线与在线两个层面消除自评乐观偏差。

\begin{figure}[htbp]
  \centering
  \includegraphics[width=\linewidth]{fig_overview}
  \caption{基于数据标定留出想象海洋的自主水下航行器计算与观测权衡方法总览。（A）海洋大数据：利用2018至2026年东海3197天逐日$2\text{ km}$卫星叶绿素与海表温度数据切出422个赤潮窗口（开发集274个、测试集148个），标定1至3天滞后底图误差（0.125、0.176、0.210）与四季演变图谱，而550万参数的深度时空预测模型并不优于直接沿用最新底图，表明小尺度结构必须靠现场观测获取。（B）出发前：围绕3天前旧图生成48个经数据标定不确定性的想象海洋，划分为32个搜索用海洋与16个留出海洋；下水前以4个留出海洋为假想真值完整预演五种用时策略（柱状图为2025年4月4日测试任务在中等计算机上的假想遗憾，“一直走”以0.715入选），选出平均假想遗憾最低的策略$s^*$投入实海。（C）任务执行：每20分钟在走、深想与停下之间行动并随时响应最浓处查询，插图显示随规划预算增加规划器自评得分持续攀升，而留出海洋评判与真实航线价值一同趋于平缓。}
  \label{fig:overview}
\end{figure}

本文主要贡献如下：
\begin{itemize}
  \item \textbf{提出独立于搜索过程的计算价值留出评判准则：}揭示了未知海洋巡测中过度计算的根源在于规划器对搜索用海洋的过拟合，提出将卫星数据标定的想象海洋划分为搜索用海洋与留出海洋，以留出海洋上的结果作为衡量计算收益的无偏标尺。
  \item \textbf{构建硬件自适应的出发前预演与在线留出估值方法：}设计了结合局地化集合同化与交叉熵规划的五种用时策略，提出在母船或岸基计算机上以留出海洋为假想真值开展全任务预演，并将留出评判嵌入在线单步元推理，无需历史任务即可适配不同算力与功耗。
  \item \textbf{提供基于八年东海赤潮窗口与双真值模型的实证及边界刻画：}在422个赤潮窗口、四种硬件与两种真值模型上开展配对检验，表明深度时空模型无法替代现场采样、留出评判将跨三个半数量级推演预算的价值偏差控制在0.005以内、停下复测在全参数空间均得不偿失，且出发前预演与在线留出估值均达到依赖历史任务方法的水平。
\end{itemize}

本文其余章节安排如下：第2节回顾相关工作；第3节建立巡测模型并剖析问题矛盾；第4节阐述留出想象海洋、出发前预演与在线留出估值方法；第5节报告实验结果与机理分析；第6节总结全文。
